% Gesamttest «Temperatur», Fassung 1.2 — Quelle des PDFs (python3 scripts/build-lp-pdf.py temperatur).
% Musterlösung und Raster: bewertungspaket.tex. Alle Werte mit python3 nachgerechnet (08.10.2026, Fassung 1.2 nach der zweiten /lp-pruefung).
% Jede Aufgabe fragt anders als Aufgaben, Übungen, Kontrollfragen, Leisten, Clipbeispiele und die Themenseite:
% G1 zwei Spuren am Bild deuten und die Temperatur als Mittelwert erklären (Kapitel 1), G2 Stoff aus Schmelz- und
% Siedepunkten rückwärts bestimmen, Kräfte gegen Bewegung, Verdunsten (Kapitel 2), G3 eigene Messreihe eintragen und
% verlängern, d) mit den Messwerten prüfen, zu welcher Skala der Druck proportional ist (Kapitel 3; anders als
% Aufgabe 3d, die Anwendungen aufzählen lässt), G4 vier Aussagen beurteilen, eine davon richtig (Kapitel 3),
% G5 Satellit in beiden Skalen (Kapitel 4), G6 Verdoppeln der Bewegungsenergie rückwärts (Kapitel 4, 3).
% Bekannte Nähe: G3 a–c hat den Aufbau von Aufgabe 3b und der Übung «nullpunkt» (dort ohne Zeichnen); bewusst belassen,
% weil das Verlängern einer eigenen Messreihe der Kern von K2 «Ursprung» ist.
% Die Spuren in G1 sind schematische Zufallswege (python3 random.Random, Spur 1: Seed 8, Schritt 0.26 cm;
% Spur 2: Seed 123, Schritt 0.42 cm; je 22 Schritte; Faktor 1.6 statt 2.5, siehe Prüfung 07.10.2026).
\documentclass[11pt]{article}
\usepackage{../lp-druck}
\renewcommand{\lpname}{Temperatur}
\renewcommand{\lpteilgebiet}{5.1}
\renewcommand{\lpfuss}{Gesamttest · 25 Punkte}
\begin{document}
\lpkopf{Gesamttest}{%
  \vspace{4pt}\noindent\begin{tabularx}{\linewidth}{@{}X X@{}}
  Code (kein Name): \hrulefill & Punkte: \hrulefill\ / 25
  \end{tabularx}}

\begin{lpinfo}{Anleitung}
\textbf{Zeit:} rund 30 Minuten. \textbf{Hilfsmittel:} Taschenrechner und Formelsammlung, in allen Teilen.
Schreib zu jeder Rechnung zuerst die Formel, dann die Werte \textbf{mit Einheiten}. Punkte gibt es für den Weg,
für das Ergebnis mit Einheit, für das Ablesen und Zeichnen im Diagramm und für Begründungen in eigenen Worten.
Umrechnen: $T\,[\text{K}] = \vartheta\,[^\circ\text{C}] + 273.15$.
Danach: Lösung scannen oder fotografieren (ohne Namen und Standort) und mit dem \emph{Bewertungspaket} bewerten lassen.
\end{lpinfo}

\lpteil{Teil A · Temperatur und Aggregatzustände}{Kapitel 1 und 2 · 8 Punkte}

\begin{lpaufgabe}{G1}{Zwei Spuren unter dem Mikroskop}{3 P}
Eine Schülerin filmt dasselbe Fetttröpfchen in Wasser zweimal je eine Minute lang: einmal bei $10\,^\circ\text{C}$, einmal bei $50\,^\circ\text{C}$. Das Bild zeigt die beiden Wege des Tröpfchens schematisch im gleichen Massstab.\par\smallskip
\begin{center}\begin{tikzpicture}[line join=round]
\draw[lplinie] (-2.6,-1.9) rectangle (2.6,1.9); \draw[lplinie] (3.2,-1.9) rectangle (8.4,1.9);
\draw[thick] (-0.48,-0.88) -- (-0.44,-0.63) -- (-0.19,-0.69) -- (-0.00,-0.50) -- (-0.08,-0.75) -- (0.15,-0.62) -- (0.15,-0.36) -- (0.41,-0.36) -- (0.48,-0.11) -- (0.31,-0.31) -- (0.06,-0.25) -- (-0.19,-0.17) -- (-0.45,-0.16) -- (-0.35,0.08) -- (-0.23,-0.15) -- (-0.01,-0.01) -- (0.02,0.25) -- (0.28,0.28) -- (0.25,0.54) -- (0.03,0.68) -- (0.24,0.53) -- (0.05,0.71) -- (0.25,0.88);
\begin{scope}[shift={(5.8,0)}]\draw[thick] (-0.74,-1.49) -- (-0.35,-1.35) -- (0.01,-1.13) -- (-0.34,-0.90) -- (-0.01,-0.64) -- (0.33,-0.88) -- (0.74,-0.78) -- (0.33,-0.88) -- (0.12,-0.51) -- (0.37,-0.85) -- (0.60,-0.50) -- (0.38,-0.14) -- (0.17,0.23) -- (0.18,0.64) -- (0.60,0.65) -- (0.22,0.81) -- (0.57,1.03) -- (0.23,0.79) -- (0.61,0.97) -- (0.44,1.35) -- (0.04,1.49) -- (0.39,1.25) -- (0.74,1.48);\end{scope}
\node at (0,-2.25) {\small Spur 1}; \node at (5.8,-2.25) {\small Spur 2};
\end{tikzpicture}\end{center}
\begin{enumerate}[label=\alph*),leftmargin=*,itemsep=2pt,topsep=2pt]
\item Welche Spur entstand bei $50\,^\circ\text{C}$?
\item Woran erkennst du es im Bild? Erkläre mit dem Teilchenmodell, warum die Spur bei höherer Temperatur so aussieht.
\item Wofür ist die Temperatur des Wassers ein Mass? Erkläre damit, warum ein einzelnes Wasserteilchen keine Temperatur hat.
\end{enumerate}
\lpplatz{6}
\end{lpaufgabe}

\begin{lpaufgabe}{G2}{Welcher Stoff?}{5 P}
Schmelz- und Siedepunkte bei Normaldruck:\par\smallskip
\begin{center}\begin{tabular}{lrrrrr}\toprule
 & Wasser & Ethanol & Quecksilber & Brom & Essigsäure\\\midrule
Schmelzpunkt & $0\,^\circ\text{C}$ & $-114\,^\circ\text{C}$ & $-39\,^\circ\text{C}$ & $-7\,^\circ\text{C}$ & $17\,^\circ\text{C}$\\
Siedepunkt & $100\,^\circ\text{C}$ & $78\,^\circ\text{C}$ & $357\,^\circ\text{C}$ & $59\,^\circ\text{C}$ & $118\,^\circ\text{C}$\\\bottomrule
\end{tabular}\end{center}
\begin{enumerate}[label=\alph*),leftmargin=*,itemsep=2pt,topsep=2pt]
\item Welcher Stoff ist bei $-20\,^\circ\text{C}$ fest und bei $70\,^\circ\text{C}$ gasförmig?
\item In welchem Temperaturbereich sind Quecksilber und Ethanol flüssig, Wasser aber fest?
\item Der Stoff aus a) wird von $-20\,^\circ\text{C}$ auf $70\,^\circ\text{C}$ erwärmt. Nenne die beiden Übergänge mit ihrer Temperatur. Begründe mit dem Teilchenmodell, warum der Stoff bei $-20\,^\circ\text{C}$ fest und bei $70\,^\circ\text{C}$ gasförmig ist.
\item Ethanol steht bei Zimmertemperatur in einem offenen Schälchen. Nach einem Tag ist es verschwunden, obwohl es erst bei $78\,^\circ\text{C}$ siedet. Begründe mit dem Teilchenmodell.
\end{enumerate}
\lpplatz{9}
\end{lpaufgabe}

\clearpage
\lpteil{Teil B · Celsius und Kelvin}{Kapitel 3 · 9 Punkte}

\begin{lpaufgabe}{G3}{Eine eigene Messreihe}{5 P}
Ein Gasthermometer (Glaskolben mit Luft, Volumen fest) wird in drei Bäder gestellt:\par\smallskip
\begin{center}\begin{tabular}{lrrr}\toprule
Bad $\vartheta$ & $-20\,^\circ\text{C}$ & $40\,^\circ\text{C}$ & $90\,^\circ\text{C}$\\
Druck $p$ & $834\;\text{hPa}$ & $1032\;\text{hPa}$ & $1196\;\text{hPa}$\\\bottomrule
\end{tabular}\end{center}
\begin{center}\begin{tikzpicture}[x=0.032cm,y=0.0075cm]
\foreach \x in {-300,-275,...,100} \draw[lplinie!70,very thin] (\x,0) -- (\x,1400);
\foreach \y in {0,50,...,1400} \draw[lplinie!70,very thin] (-300,\y) -- (100,\y);
\foreach \x in {-300,-200,...,100} {\draw[lplinie] (\x,0) -- (\x,1400); \node[below] at (\x,0) {\small $\x$};}
\foreach \y in {200,400,...,1400} {\draw[lplinie] (-300,\y) -- (100,\y); \node[left] at (-300,\y) {\small \y};}
\node[left] at (-300,0) {\small 0};
\draw[->] (-300,0) -- (112,0) node[above left] {\small $\vartheta$ in $^\circ$C};
\draw[->] (-300,0) -- (-300,1460) node[right] {\small $p$ in hPa};
\end{tikzpicture}\end{center}
\begin{enumerate}[label=\alph*),leftmargin=*,itemsep=2pt,topsep=2pt]
\item Trage die drei Messpunkte ins Diagramm ein und zeichne die Gerade bis zum Druck null.
\item Lies ab, bei welcher Temperatur deine Gerade den Druck null erreicht.
\item Berechne diese Temperatur aus den Messpunkten bei $40\,^\circ\text{C}$ und $90\,^\circ\text{C}$.
\item Prüfe mit den Messpunkten bei $40\,^\circ\text{C}$ und $90\,^\circ\text{C}$: Ist der Druck proportional zur Celsius-Temperatur oder zur Kelvin-Temperatur? Was folgt daraus, wann man in Kelvin rechnen muss und wann Celsius genügt?
\end{enumerate}
\lpplatz{6}
\end{lpaufgabe}

\begin{lpaufgabe}{G4}{Stimmt das?}{4 P}
Beurteile jede Aussage mit richtig oder falsch, und begründe in einem Satz.
\begin{enumerate}[label=\alph*),leftmargin=*,itemsep=2pt,topsep=2pt]
\item «Am Morgen waren es $5\,^\circ\text{C}$, am Mittag $10\,^\circ\text{C}$: Die Luftteilchen haben jetzt die doppelte mittlere Bewegungsenergie.»
\item «Eis schmilzt bei $0\,^\circ\text{C}$, Wasser siedet bei $100\,^\circ\text{C}$. Darum siedet Wasser auch auf dem Jungfraujoch bei $100\,^\circ\text{C}$.»
\item «Ein Thermometer im Labor zeigt $-5\;\text{K}$.»
\item «Wird Wasser um $15\;\text{K}$ erwärmt, ist das eine Erwärmung um $15\,^\circ\text{C}$.»
\end{enumerate}
\lpplatz{7}
\end{lpaufgabe}

\lpteil{Teil C · Umrechnen}{Kapitel 4 · 8 Punkte}

\begin{lpaufgabe}{G5}{Ein Satellit}{4 P}
Ein Satellit hat auf der Sonnenseite eine Oberflächentemperatur von $120\,^\circ\text{C}$, auf der Schattenseite $-160\,^\circ\text{C}$.
\begin{enumerate}[label=\alph*),leftmargin=*,itemsep=2pt,topsep=2pt]
\item Gib beide Temperaturen in Kelvin an.
\item Wie gross ist der Temperaturunterschied zwischen den beiden Seiten in Kelvin?
\item Ein Sensor im Innern muss auf $40\;\text{K}$ gekühlt werden. Wie viel Grad Celsius sind das?
\end{enumerate}
\lpplatz{7}
\end{lpaufgabe}

\begin{lpaufgabe}{G6}{Doppelt so viel Bewegungsenergie}{4 P}
Luft in einer verschlossenen Stahlflasche (Volumen fest) hat $20\,^\circ\text{C}$.
\begin{enumerate}[label=\alph*),leftmargin=*,itemsep=2pt,topsep=2pt]
\item Auf welche Temperatur in Grad Celsius muss man die Flasche erwärmen, damit die mittlere Bewegungsenergie der Luftteilchen doppelt so gross wird?
\item Um welchen Faktor steigt dabei der Druck in der Flasche? Begründe.
\item Um wie viel Kelvin hat sich die Temperatur in a) erhöht?
\end{enumerate}
\lpplatz{7}
\end{lpaufgabe}

\vfill
{\small\color{lpgrau}Bewerten: mit dem Bewertungspaket (Musterlösung, Punkteraster, Auftrag an eine KI) — erst nach dem Lösen öffnen.}
\end{document}
